\documentclass[11pt]{article}
\usepackage[a4paper,top=0.6cm,bottom=3.0cm,left=1.9cm,right=1.9cm,footskip=2.0cm]{geometry}
\usepackage{graphicx}
\usepackage[table]{xcolor}
\usepackage{ragged2e}
\usepackage{fancyhdr}
\usepackage{fontspec}
\setmainfont{Noto Sans Gurmukhi}
\setsansfont{Noto Sans Gurmukhi}
\newfontfamily\latinfont{Carlito}
\renewcommand{\familydefault}{\sfdefault}
\definecolor{titleblue}{HTML}{0F2F4A}
\definecolor{textblack}{HTML}{1A1A1A}
\definecolor{rule}{HTML}{C9C4C0}
\definecolor{muted}{HTML}{6B635E}
\setlength{\parindent}{0pt}\setlength{\parskip}{0.5em}

\setlength{\headheight}{28pt}\setlength{\headsep}{6pt}
\newcommand{\idmfooter}{%
  \begin{minipage}{\textwidth}
    {\color{rule}\hrule height 0.7pt}\vspace{4pt}
    \noindent
    \begin{minipage}[c]{0.66\textwidth}{\scriptsize\color{muted}{\latinfont Developed by the Water and Climate Lab, IIT Gandhinagar, using hydro-meteorological data from partner agencies}.}\end{minipage}\hfill
    \begin{minipage}[c]{0.32\textwidth}\raggedleft
      
\includegraphics[height=0.62cm]{wcl.png}\hspace{6pt}%
      
\includegraphics[height=0.62cm]{iitgn.png}
    \end{minipage}\\[3pt]
    {\scriptsize\color{muted}© {\latinfont 2026 Water and Climate Lab} · {\latinfont Indian Institute of Technology Gandhinagar} · {\latinfont For research and demonstration purposes}.}
  \end{minipage}%
}
\newcommand{\idmrunhead}{%
  \begin{minipage}{\textwidth}
    {\small\bfseries\color{titleblue}{\latinfont National Drought Summary}}\hfill{\small\color{muted}{\latinfont Week ending June 3, 2026}}\\[2pt]
    {\color{rule}\hrule height 0.6pt}
  \end{minipage}%
}
\fancypagestyle{idmfoot}{\fancyhf{}\renewcommand{\headrulewidth}{0pt}\renewcommand{\footrulewidth}{0pt}\fancyfoot[C]{\idmfooter}}
\fancypagestyle{idmrun}{\fancyhf{}\renewcommand{\headrulewidth}{0pt}\renewcommand{\footrulewidth}{0pt}\fancyhead[C]{\idmrunhead}\fancyfoot[C]{\idmfooter}}
\pagestyle{idmrun}

\begin{document}

\thispagestyle{idmfoot}
\noindent
\begin{minipage}[c]{0.16\textwidth}
\includegraphics[width=\linewidth]{wcl.png}\end{minipage}\hfill
\begin{minipage}[c]{0.80\textwidth}\raggedleft
  {\fontsize{19}{22}\selectfont\bfseries\color{titleblue}{\latinfont National Drought Summary}}\\[2pt]
  {\normalsize\color{muted}{\latinfont Week ending June 3, 2026}}
\end{minipage}
\vspace{6pt}{\color{rule}\hrule height 1pt}\vspace{10pt}
\begin{center}
  \includegraphics[width=0.60\textwidth]{cdi_map.png}\\[3pt]
  {\small\color{muted}{\latinfont Combined Drought Index (CDI) for the week ending June 3, 2026}.}
\end{center}
\vspace{2pt}
{\large\bfseries\color{titleblue}{\latinfont Summary}}\\[3pt]
{\justifying\normalsize\color{textblack}ਇਸ ਹਫ਼ਤੇ ਦੇ ਰਾਸ਼ਟਰੀ ਬੇਅ ਦੀ ਸਥਿਤੀ ਵਿੱਚ ਬੇਅ ਸ਼੍ਰੇਣੀਆਂ ਵਿੱਚ ਵੰਡ ਦਿਖਾਈ ਦਿੰਦੀ ਹੈ ਜਿਸ ਵਿੱਚ {\latinfont 80.4}\% ਖੇਤਰ ਆਮ ਸ਼੍ਰੇਣੀ ਵਿੱਚ ਅਤੇ {\latinfont 19.6}\% ਡੀ{\latinfont 0}+ ਸ਼੍ਰੇਣੀ ਵਿੱਚ ਹੈ। ਇਸ ਹਫ਼ਤੇ ਡੇਅ ਨਾਲ ਪ੍ਰਭਾਵਿਤ ਕੁੱਲ ਖੇਤਰ {\latinfont 19.6}\% ਸੀ, ਜੋ ਪਿਛਲੇ ਹਫ਼ਤੇ ਦੇ {\latinfont 20.1}\% ਨਾਲੋਂ {\latinfont 0.5} ਪ੍ਰਤੀਸ਼ਤ ਬਿੰਦੂ ਦਾ ਸੰਕੋਚ ਹੈ। ਪਿਛਲੇ ਮਹੀਨੇ ਅਪ{\latinfont asan 15.0}\% ਖੇਤਰ ਡੀ{\latinfont 0}+ ਵਿੱਚ ਸੀ। ਗੁਜਰਾਤ ਨੇ ਸਭ ਤੋਂ ਵੱਧ ਬੇਅ ਪ੍ਰਭਾਵ ਦਿਖਾਇਆ, ਜਿਸਦਾ {\latinfont 85.2}\% ਖੇਤਰ ਬੇਅ ਵਿੱਚ ਹੈ। ਦਿੱਲੀ ਕੇਂਦਰ ਸ਼ਾਸਤ ਪ੍ਰਦੇਸ਼ ਵਿੱਚ {\latinfont 69.6}\%, ਅਰੁਣਾਚਲ ਪ੍ਰਦੇਸ਼ ਵਿੱਚ {\latinfont 52.8}\%, ਜਮੁਕਸ਼ਮ{\latinfont IR} ਵਿੱਚ {\latinfont 50.4}\%, ਮਿਜ਼ੋਰਮ ਵਿੱਚ {\latinfont 49.8}\%, ਅਤੇ ਮੇਘਾਲਿਆ ਵਿੱਚ {\latinfont 43.5}\% ਬੇਅ ਦੀ ਮੌਜੂਦਗੀ ਹੈ। ਇਸ ਦੇ ਉਲਟ, ਸਿੱਕਿਮ ਨੇ {\latinfont 0.0}\% ਅਤੇ ਨਗਾਲੈਂਡ ਨੇ {\latinfont 0.7}\% ਬੇਅ ਦਰ ਰਿਪੋਰਟ ਕੀਤੀ।\par}
\vspace{8pt}
{\large\bfseries\color{titleblue}{\latinfont Regional Outlook}}\\[2pt]
{\normalsize\color{titleblue}\textbf{{\latinfont North}}}\enspace {\normalsize\color{textblack}{\latinfont Moderate, scattered drought. About 30}\% {\latinfont of the area is in drought (D0}–{\latinfont D4). Pockets of severe-or-worse conditions are present}.} \par\smallskip {\normalsize\color{titleblue}\textbf{{\latinfont Northwest}}}\enspace {\normalsize\color{textblack}{\latinfont Mostly normal, with localised dryness. About 5}\% {\latinfont of the area is in drought (D0}–{\latinfont D4)}.} \par\smallskip {\normalsize\color{titleblue}\textbf{{\latinfont Northeast}}}\enspace {\normalsize\color{textblack}{\latinfont Moderate, scattered drought. About 37}\% {\latinfont of the area is in drought (D0}–{\latinfont D4). A notable share is in severe-or-worse categories (D2}–{\latinfont D4)}.} \par\smallskip {\normalsize\color{titleblue}\textbf{{\latinfont East}}}\enspace {\normalsize\color{textblack}{\latinfont Mostly normal, with localised dryness. About 11}\% {\latinfont of the area is in drought (D0}–{\latinfont D4). Pockets of severe-or-worse conditions are present}.} \par\smallskip {\normalsize\color{titleblue}\textbf{{\latinfont Central}}}\enspace {\normalsize\color{textblack}{\latinfont Mostly normal, with localised dryness. About 11}\% {\latinfont of the area is in drought (D0}–{\latinfont D4). Pockets of severe-or-worse conditions are present}.} \par\smallskip {\normalsize\color{titleblue}\textbf{{\latinfont West}}}\enspace {\normalsize\color{textblack}{\latinfont Widespread drought. About 42}\% {\latinfont of the area is in drought (D0}–{\latinfont D4). Pockets of severe-or-worse conditions are present}.} \par\smallskip {\normalsize\color{titleblue}\textbf{{\latinfont South}}}\enspace {\normalsize\color{textblack}{\latinfont Mostly normal, with localised dryness. About 8}\% {\latinfont of the area is in drought (D0}–{\latinfont D4)}.}
\end{document}
